% Latin 17859, batch 3 — Gallica views 41–49.
% Pass: Opus 5 (claude-opus-5), 2026-09-15 — first pass, unchecked against
% the views by a human.
%
% What the views are. Every view is a microfilm image (greyscale, landscape,
% about 7300 × 5850) of the volume lying open: a left page and a right page
% side by side, the gutter near the middle. Each view was split into its two
% halves and read as two pages; in the file each \page{N} holds the left page
% then the right, and a \note{} says where the right page begins. The volume
% ends inside this batch: view 49 is the last view of Latin 17859, and its
% right page is blank. No view was skipped; 49 right is the only blank page.
%
% What the nine views hold. A copy, in one regular italic hand throughout,
% of Maurolico's spherical trigonometry and of what he appended to it.
%
%   v. 41 L – 42 L, 45 R – 47 L  The end of a book of propositions on
%             triangles « ex arcubus circulorum majorum in superficie
%             sphæræ »: the end of Prop. XXIII (continuing view 40), then
%             XXIIII–XXXII, each with its enunciation, demonstration and
%             marginal figures, two « Corollaria », and a « Scholium » that
%             turns to worked examples on declinations and ascensions. The
%             demonstrations cite « præmissi libelli », « præcedentis libri »,
%             Euclid's Elements, Theodosius's and Menelaus's Sphaerica.
%   v. 42 R – 45 L  A continuous text on the tables of « Iohannes de
%             Monteregio » (Regiomontanus): why his rules for declination,
%             right ascension and ascensional difference work (the division
%             by the sinus totus done by striking five figures, the unit
%             added when the struck figures exceed 50000), then the
%             « Demonstratio tabulæ beneficæ » — a table made « ad
%             imitationem tabulæ foecundæ », with its use.
%   v. 47 R   The end of that text, dated in the first person: « Hæc in Arce
%             Apollinari dum cum Iohanne Vigintimillio Hieracensium
%             Marchione degeremus olim mense Augusto 1530 speculabamur ».
%             Then a dedicatory letter « D. Octavio Spinulæ Quæstori ærario
%             Siciliæ, Regioque Consiliario, Maurolycus. Salutem. », a
%             « Hexastichon » to the same, and the date « Messanæ. Kal. Junij
%             M.D.LVII. ».
%   v. 48 L   A new text: « Autolyci de sphæra quæ movetur ex traditione
%             Maurolyci. Liber. », the definition and Props. 1–4, breaking
%             off at the foot of the page.
%   v. 48 R – 49 L  Numbered problems (« 7 », « 8 ») worked with the tabula
%             benefica and full numerical examples (Sun in 11° Taurus,
%             latitude 38°), ending with a rule for rounding. Then the line
%             « Compendium Mathematicæ. omittitur quia a P. M. Mersenno
%             editum in Apologia scientiarum cap. ultimo. », and nothing
%             after it. That line is the copyist's, and it is the one
%             internal sign in the batch of when the copy was made: after
%             a publication by Mersenne. This is recorded as read; the
%             catalogue's dating, 1601-1700, is copied as it stands.
%
% The order of the leaves. Read view by view, the text runs from each right
% page onto the next view's left page without a break (the two sides of one
% leaf). Inside an opening it runs on at views 41, 43, 44 and 46, and breaks
% at views 42, 45, 47 and 48:
%   – 42 L stops inside Prop. XXVIII (« quadrangulum quod »), and 45 R
%     repeats that page's last two lines and finishes the proposition;
%   – 45 L stops on « vel quot zyphræ sunt in numero », and 47 R opens on
%     « sinus totius, Et habeo in quotiente sinum arcus cd quæsiti »;
%   – 47 L stops on « Si autem sumatur arcus fbc pro horizonte recto », and
%     42 R opens on « erit arcus ac ascensio recta, respondens arcui
%     Zodiaci ab ».
% So the three leaves inked 41, 42, 43 (42 R – 45 L) read as if they belonged
% between the leaves inked 45 and 46, i.e. after 47 L and before 47 R. That
% is this pass's reconstruction from the sense; the ink numbers run 40–47
% without a gap in the present order, and the file keeps the present order.
% 48 L (Autolycus) does not continue 47 R (which ends on a date) and is not
% continued by 48 R: it may be the beginning of a text whose leaves are
% elsewhere, or a text abandoned after one page. Not settled.
%
% The numbers on the leaves. Each right page from view 41 to view 48 carries
% at the top right a number in ink, larger than the text: 40, 41, 42, 43,
% 44, 45, 46, 47. One number per leaf, on the right page only, consecutive.
% They are ink, not the thin pencil of a library foliation, and the stroke
% could be the copyist's; this pass cannot tell whether they are the
% foliation the literature cites for Latin 17859. No \folio{} is written;
% each number is recorded in the \note{} opening its page. Left pages carry
% no number; 49 R carries none. Other numbers: « 7 » over a line and « 8. »
% in the margin of 48 R (problem numbers); proposition numbers in the margin
% of 48 L. A round stamp « Bibliothèque impériale MSS » in the margin of
% 49 L.
%
% Letters, glyphs, notation. Figure letters are set in math ($abc$), as the
% hand writes them run together (« arcus ab », « triangulum abc »). One
% figure letter is a reversed c, which the PDF fonts have no glyph for; it is
% set $\supset$, the nearest standard symbol (view 46 L, in the text and in
% the figure note). It is a letter, not inclusion.
% The hand uses the long s, set s, and one looped initial form for u and v,
% set u or v by sound; its initial i, which looks like a capital J
% (« In », « Iam », « Igitur »), is set lower case inside a sentence. « æ »
% is written as a hooked e and set æ. The capital « Et » the hand writes
% inside sentences is kept where it is clear. Contractions are kept as
% written — « dria », « driam », « driæ », « diff\textsuperscript{ia} »,
% « quæstionib\textsuperscript{9} », « nram », « comp\textsuperscript{tum} »,
% « propositionẽ » — and expanded in a \note{} at first occurrence. The
% digit 0 is often an « a »-like loop; « 11 » inside a number is often a
% « u » (« 2u89 » = 21189, checked against 61389 + 21189 = 82578 on
% view 49). Numbers are copied as written, errors included (views 48 R,
% 49 L). The crescent sign « ( » that opens several paragraphs (43, 44, 45,
% 47) is the hand's paragraph mark; it is not transcribed, and a \note{}
% records where it stands. Figures are not drawn; each page's figures are
% described, with their lettering, in a \note{}. Headings (« Propositio
% XXV », « Corollaria », « Scholium », the letter's address) are set as
% \subsection*; the Autolycus title as \section*.
%
% Illegible: 3. Uncertain readings: 13.
%
% Nothing here was checked against any printed edition of Maurolico or any
% other transcription.
\documentclass[11pt,a4paper]{article}
% The preamble shared by every transcription of the mathematicians' manuscripts in Gallica.
%
% It rests on one idea: **uncertainty compounds, it does not resolve.** A
% transcription that smooths over a doubtful word has destroyed the only thing
% it was made for — a reader must be able to tell, at every point, what was on
% the page and what was supplied. The apparatus macros below are therefore the
% critical apparatus, not decoration.
%
% The same commands are recognised by `scripts/render.mjs`, which produces the
% site's reading view, and by `scripts/tei.mjs`. A macro added here must be
% added there too, or rendering fails loudly — which is the wanted behaviour.
%
% Comments here are English, like the rest of the repository. The documents
% this preamble sets are French, and that is deliberate: see the skills.

\ProvidesPackage{germain}[2026/09/15 Transcriptions of mathematicians' manuscripts in Gallica]

\usepackage{amsmath,amssymb,amsthm}
\usepackage{mathrsfs}
% Kept from the parent preamble so the renderer's diagram support stays
% available; these manuscripts draw no commutative diagrams, and nothing here uses it.
\usepackage{tikz-cd}
\usepackage[english,french]{babel}
\usepackage{fontspec}

% --- Greek ------------------------------------------------------------------
% The text face stays fontspec's default, Latin Modern, which has no Greek:
% XeTeX drops every glyph it lacks with a line in the log and nothing on the
% page, and the Greek words the manuscripts quote vanished from the PDFs. So
% the Greek and Greek Extended blocks get a XeTeX character class of their own,
% and entering or leaving a run of it switches the family to CMU Serif — the
% same Computer Modern design, with polytonic Greek — and back. Only the
% family changes: a Greek word in \textit or \textbf keeps its shape and series.
% The transitions to and from every other class carry the tokens that class
% has towards an ordinary letter, so French spacing before « ; » or inside
% « » still holds after a Greek word. No grouping, so no brace or \endgroup
% can come out unbalanced; maths is left alone, since XeTeX inserts nothing
% there. Kept identical in `exercices.sty`.
\makeatletter
\newfontfamily\germainGreekfont{cmun}[
  Extension=.otf, NFSSFamily=germaingreek,
  UprightFont=*rm, ItalicFont=*ti, SlantedFont=*sl,
  BoldFont=*bx, BoldItalicFont=*bi, BoldSlantedFont=*bl]
\newXeTeXintercharclass\germain@greek
\def\germain@greekrange#1#2{%
  \@tempcnta=#1\relax
  \loop
    \XeTeXcharclass\@tempcnta=\germain@greek
    \ifnum\@tempcnta<#2\advance\@tempcnta\@ne\repeat}
\germain@greekrange{"0370}{"03FF}
\germain@greekrange{"1F00}{"1FFF}
\def\germain@greekprev{\rmdefault}
\def\germain@greekin{%
  \edef\germain@greekprev{\f@family}\fontfamily{germaingreek}\selectfont}
\def\germain@greekout{\fontfamily{\germain@greekprev}\selectfont}
\def\germain@greekjoin{%
  \XeTeXinterchartokenstate=\@ne
  \@tempcnta=\z@
  \loop
    \ifnum\@tempcnta=\germain@greek\else
      \XeTeXinterchartoks\germain@greek\@tempcnta=\expandafter{%
        \expandafter\germain@greekout\the\XeTeXinterchartoks\z@\@tempcnta}%
      \XeTeXinterchartoks\@tempcnta\germain@greek=\expandafter{%
        \the\XeTeXinterchartoks\@tempcnta\z@\germain@greekin}%
    \fi
    \ifnum\@tempcnta<4095\advance\@tempcnta\@ne\repeat}
% babel-french sets its own classes, and saves and restores the interchar
% state, each time a language is selected: join again after it does.
\addto\extrasfrench{\germain@greekjoin}
\addto\extrasenglish{\germain@greekjoin}
\AtBeginDocument{\germain@greekjoin}
\makeatother

\usepackage{geometry}
\geometry{a4paper,margin=25mm,marginparwidth=28mm}
\usepackage{marginnote}
\usepackage{xcolor}
\usepackage[normalem]{ulem}
\setlength{\emergencystretch}{3em}
\usepackage{hyperref}
\hypersetup{colorlinks=true,linkcolor=black,urlcolor=black,pdfborder={0 0 0}}

% --- Batch metadata ---------------------------------------------------------
% Taken as they stand by the rendering script, which shows them at the head of
% the reading view. They fix what the file claims to cover. The macro names are
% the parent project's, so that the scripts stay shared; \folder here means the
% volume's slug — `fr-9115` — and \pages counts Gallica views.

\newcommand{\folder}[1]{\gdef\grothFolder{#1}}
\newcommand{\batch}[1]{\gdef\grothBatch{#1}}
\newcommand{\pages}[2]{\gdef\grothFirst{#1}\gdef\grothLast{#2}}
\newcommand{\foldertitle}[1]{\gdef\grothFoldertitle{#1}}

% The holder's own shelfmark, printed as they write it: « Français 9115 ».
\newcommand{\shelfmark}[1]{\gdef\grothShelfmark{#1}}

% The dating, copied verbatim from the holder's catalogue — never inferred
% here. « XIXe siècle » is the BnF's; a pass that dates a leaf from its content
% says so in a \note{}, as its own inference.
\newcommand{\dating}[1]{\gdef\grothDating{#1}}

% The status notice, printed in the title block and repeated as a diagonal
% watermark on every page of the PDF. The manuscripts are in the public
% domain; what this line declares is not a right but a quality — an
% unverified first pass by a machine. The skills write the line; a file
% without it gets no watermark, so leaving it out is visible in the output.
\newcommand{\watermark}[1]{\gdef\grothWatermark{#1}}

\gdef\grothFolder{?}\gdef\grothBatch{}\gdef\grothFirst{?}\gdef\grothLast{?}%
\gdef\grothFoldertitle{}\gdef\grothShelfmark{}\gdef\grothDating{}\gdef\grothWatermark{}

% --- The résumé -------------------------------------------------------------
\newenvironment{resume}
  {\small\setlength{\parskip}{0.45em}}
  {\par\medskip\hrule\medskip}

% --- Keywords ---------------------------------------------------------------
% In English, closing the modernised reading's résumé: the modern vocabulary
% under which to search for what the leaves construct. The manifest extracts
% them to tag the volume — this line is the single source of the tags.
\newcommand{\keywords}[1]{\par\smallskip\noindent
  {\footnotesize\textsc{Keywords} --- \textit{#1}}\par}

% --- The critical apparatus -------------------------------------------------

% The Gallica view number, in the margin. This is the anchor that lets a
% disagreement over a formula be settled by looking at *one* image rather than
% a volume of 750. The site uses it to turn the facsimile as the transcription
% is scrolled. A view is one scanned image: usually one side of a leaf.
\newcommand{\page}[1]{%
  \marginnote{\footnotesize\color{gray!70}\textsf{v.\,#1}}%
  \label{page:#1}\ignorespaces}

% The BnF's pencilled foliation, where the leaf carries it and a pass can read
% it: « 198r ». This is what the literature cites — « ff. 198r–208v » — and
% the one line that turns a citation into a view. Recorded, never inferred:
% a folio a pass cannot read is not written.
\newcommand{\folio}[1]{%
  \marginnote{\footnotesize\color{gray!50}\textsf{f.\,#1}}\ignorespaces}

% The modernised reading's coarser counterpart to \page{}: which views a
% section is drawn from.
\newcommand{\pagerange}[2]{%
  \marginnote{\footnotesize\color{gray!70}\textsf{v.\,#1--#2}}%
  \ignorespaces}

% Illegible: never guessed.
\newcommand{\ill}{\textcolor{red!60!black}{[\,\ldots\,]}}

% A doubtful reading: the word is given, and flagged as uncertain.
\newcommand{\uncertain}[1]{\uline{#1}}

% An editorial addition — a missing letter, an implied word.
\newcommand{\add}[1]{\textcolor{blue!50!black}{[#1]}}

% Struck out by the writer. What was crossed out is part of the document.
\newcommand{\struck}[1]{\sout{#1}}

% The transcriber's note, on the state of the page or the mathematics.
\newcommand{\note}[1]{\par\smallskip{\small\color{gray!60!black}\textit{[#1]}}\par\smallskip}

% A marginal note of hers, distinct from the body of the text.
\newcommand{\marginal}[1]{\par\smallskip\noindent\hspace{1em}%
  {\small\color{gray!60!black}#1}\par\smallskip}

% --- Title ------------------------------------------------------------------
% The title block is French, like the documents themselves.

\AddToHook{shipout/background}{%
  \ifx\grothWatermark\empty\else
    \begin{tikzpicture}[remember picture, overlay]
      \node[rotate=52, scale=3.4, align=center, text=gray!18,
            font=\sffamily\bfseries]
        at (current page.center) {\grothWatermark};
    \end{tikzpicture}%
  \fi}

\AtBeginDocument{%
  \begin{center}
    {\large\bfseries Manuscrits de mathématiciens (Gallica) ---
      \ifx\grothShelfmark\empty\grothFolder\else\grothShelfmark\fi}\\[2pt]
    {\itshape\grothFoldertitle}\\[4pt]
    {\small vues \grothFirst--\grothLast
      \ifx\grothBatch\empty\else\ (lot \grothBatch)\fi}\\[2pt]
    \ifx\grothDating\empty\else
      {\small\color{gray!60!black}Datation du catalogue : \grothDating}\\[2pt]
    \fi
    \ifx\grothWatermark\empty\else
      {\small\bfseries\color{red!45!gray}\grothWatermark}\\[2pt]
    \fi
    {\footnotesize\color{gray!60!black}Transcription établie d'après les images IIIF de Gallica.
     Source gallica.bnf.fr / Bibliothèque nationale de France.\\
     Les lectures incertaines sont soulignées, les illisibles notées [\,\ldots\,],
     les ajouts entre crochets ; \textsf{v.} numérote les vues de Gallica,
     \textsf{f.} la foliotation au crayon de la BnF.}
  \end{center}
  \vspace{1em}}

\endinput


\folder{latin-17859}
\shelfmark{Latin 17859}
\batch{3}
\pages{41}{49}
\dating{1601-1700}
\watermark{Lecture automatique\\non vérifiée}
\foldertitle{Opuscules de François Maurolyco sur les mathématiques et l'astronomie.}

\begin{document}

\page{41}
\note{la vue 41 montre une double page ouverte ; ce qui suit est la page de gauche. Le texte continue la proposition dont le début est à la vue 40 (hors de ce lot) : la dernière ligne de la vue 40 s'arrête sur « quod ».}
impossibile. Omnino igitur existente angulo $c$ recto rectus erit et
$f$ angulus. Et ob id in eo casu æquales erunt $c$ $f$ anguli. Hinc
sequitur ut si rectus non sit $c$ angulus neque iam rectus esse
possit $efd$ angulus. Tunc ergo ducatur circulorum majorum
arcus $eg$ $bh$ perpendiculares ad ipsos arcus $df$ $ac$. Qui quidem
perpendiculares aut ambo cadunt intra $abc$ $def$ triangula
aut unus intra, alter extra. Cadant primum ambo intus.
Et tunc quoniam per hypothesim et conversam proportionem
sicut sinus arcus $bc$ ad sinum arcus $ba$ sicut sinus arcus $fe$
ad sinum arcus $ed$. Propterea ex æquali sinus arcus $cb$ ad sinum
arcus $bh$, erit sicut sinus arcus $fe$ ad sinum arcus $eg$. Quare per
8 præmissi quoniam recti sunt qui ad $g$ $h$ æquales erunt
$c$ $f$ anguli sicut prius. Cadant deinde perpendiculares $eg$ $bh$
alter intus alter extra, hoc est $bh$ intra triangulum $abc$. Et
arcus $eg$ extra. Dico tunc quod angulus $c$ cum angulo $efd$
iunctus duos rectos faciet, Quod sic ostendam. Rursus enim per
hypothesim, et conversim sinus arcus $bc$ ad sinum arcus $ba$
sicut sinus arcus $fe$ ad sinum arcus $ed$. Et per 7 præmissi sinus
arcus $ab$ ad sinum arcus $bh$ sicut sinus arcus $de$ ad sinum
arcus $eg$. Quare erit et sinus arcus $cb$ ad sinum arcus $bh$ sicut
sinus arcus $fe$ ad sinum arcus $eg$. Quare ut prius per 8 præmissi
æqualis erit angulus $c$ angulo $efg$. Sed angulus $efg$ cum
angulo $efd$ conflat duos rectos. Ergo et angulus $c$ cum
angulo $efd$ conficiet duos rectos. Quod erat demonstrandum.
Atque ita constat quicquid proponitur demonstrandum.
\note{dans la marge gauche, deux figures : un triangle sphérique lettré $a$, $b$, $c$, avec l'arc $bh$ tombant sur $ac$ en $h$ ; un second lettré $e$ au sommet, $d$ à gauche et, sur la base, « $f$ », « $g$ » puis « $f$ » à l'extrémité droite — la lettre $f$ est écrite deux fois sur la figure.}

\subsection*{Propositio XXIIII}

Si ab angulo trianguli ex arcubus circulorum majorum
in superficie sphæræ constituti descendat ad basim
arcus circuli majoris angulum ipsum per æqualia dividens,
tunc sinus arcuum angulum ipsum continentes erunt
sinibus factarum basis portionum proportionales.

Sit in superficie sphæræ triangulum $abc$ ex arcubus circulorum
majorum. A cuius angulo quolibet puta $b$ descendat ad oppositum
arcum circuli majoris arcus $bd$ secans per æqualia angulum
$abc$. Dico tunc quod, sicut est sinus arcus $ab$ ad sinum arcus
$bc$, sic erit sinus arcus $ad$ ad sinum arcus $dc$. Nam cum triangulorum
sphæralium $bad$ $bcd$ anguli $abd$ $dbc$ sint per hypothesim æquales
Et anguli $adb$ $bdc$ conjuncti perficiant duos rectos. Iam per 22.
huius erit sinus arcus $ab$ ad sinum arcus $bc$ sicut sinus arcus
$ad$ ad sinum arcus $dc$. Quod est propositum.
\note{figures dans la marge : un triangle $abc$ coupé par l'arc $bd$, $d$ sur $ac$ ; plus bas, une figure lettrée $a$, $b$, $c$, $d$, $e$, $f$, qui se rapporte à la proposition suivante.}

\subsection*{Propositio XXV}

\note{page de droite de la vue 41. En haut à droite, à l'encre, le nombre « 40 ».}
Quod si sinus arcuum dictum angulum continentium
proportionales fuerint sinibus factarum basis portionum
tunc arcus descendens angulum ipsum per æqualia secare
probabitur.

Hoc est in eadem descriptione ponatur arcus $ab$ ad sinum
$bc$ sicut sinus arcus $ad$ ad sinum arcus $dc$.\note{sic : « ponatur arcus $ab$ ad sinum $bc$ », sans « sinus » devant « arcus ».} Dico tunc quod arcus
$bd$ per æqualia secat angulum $abc$. Nam si anguli ad $d$
sint recti, constat propositum per 8. præmissi libelli, Secus
autem erit angulorum ad $d$ obtusus alter, et alter acutus,
Acutus esto $bdc$. \uncertain{Producatque} arcu $bd$, et ad eum perpendiculari
arcu circuli majoris $ae$ fiat $ef$ arcus æqualis ipsi $de$. Et ducatur
arcus circuli magni $af$. Eritque per 28. 2. Sphæricorum
Elementorum arcus $af$ æqualis arcui $ad$. Et per 8. præced.
angulus $f$ æqualis angulo $ade$: eiusque contraposito $bdc$.
Et quoniam supponitur sinus arcus $ab$ ad sinum arcus $bc$
sicut sinus arcus $ad$ ad sinum arcus $dc$. Ideo erit et sinus
arcus $ab$ ad \struck{ad} sinum arcus $bc$ sicut sinus arcus $af$ ad
sinum arcus $dc$. Suntque anguli $f$, $bdc$ æquales acuti. Igitur
per 23 præcedentem. anguli oppositi arcubus $af$ $cd$ hoc est
anguli $abd$ $dbc$ aut æquales erunt aut simul juncti duobus
rectis æquales. Sed duobus rectis æquales esse non possunt,
quandoquidem angulum faciunt $ab$ $bc$ arcus. æquales ergo
erunt sicut proponitur demonstrandum. Hoc idem ab
impossibili per præcedentem demonstrari potest.

\subsection*{Propositio XXVI.}

Si bina triangula fuerint in superficie sphæræ ex arcubus
circulorum majorum, quorum unius duo anguli fuerint
duobus alterius angulis singuli singulis æquales, et ab
angulis reliquis perpendiculares arcus circulorum majorum
ad bases ducantur, Tunc sinus factarum basis portionum
in uno triangulo proportionales erunt sinibus factarum
basis portionum in reliquo triangulo.

Sint in superficie sphæræ duo triangula $abc$ $def$ ex arcubus circulorum
majorum. Quorum anguli $a$ $d$ sint æquales. Itemque anguli $c$ $f$
æquales. Arcus autem $bg$ $eh$ circulorum majorum perpendiculares
ad bases $ac$ $df$. Demonstrandum est quod sinus arcus $ag$ ad
sinum arcus $gc$ est sicut sinus arcus $dh$ ad sinum arcus $hf$ hoc modo.
Per 13. præcedentis libelli, ratio sinus arcus $ag$ ad sinum arcus
$gc$ componitur ex rationibus sinus arcus $ab$ ad sinum arcus $bc$.
Et sinus anguli $abg$ ad sinum anguli $gbc$. Et per eandem 13.
ratio sinus arcus $dh$ ad sinum arcus $hf$ componitur ex rationibus
\note{dans la marge droite, deux figures : un triangle $abc$ avec l'arc $bg$, $g$ sur $ac$ ; un triangle $def$ avec l'arc $eh$, $h$ sur $df$. Le texte de la page s'arrête ici au bas du feuillet ; il reprend en tête de la vue 42.}

\page{42}
\note{double page ; page de gauche. Le texte continue sans interruption la fin de la vue 41.}
sinus arcus $de$ ad sinum arcus $ef$. Et sinus anguli $deh$ ad sinum
anguli $hef$. Sed per 16 præmissi quoniam anguli $a$ $d$ æquales
Et anguli $c$ $f$ æquales sinus arcus $ab$ ad sinum arcus $bc$ sicut
sinus arcus $de$ ad sinum arcus $ef$. Et per 22. eiusdem, propter
dictorum angulorum æqualitatem sinus anguli $abg$ ad sinum
anguli $gbc$, sicut sinus anguli $deh$ ad sinum anguli $hef$.
Igitur ex æqua proportione sinus arcus $ag$ ad sinum arcus $gc$
erit sicut sinus arcus $dh$ ad sinum arcus $hf$. quod proponitur
demonstrandum.
\note{le chiffre lu « 16 » a la forme du 6 en « b » de cette main.}

\subsection*{Propositio XXVII.}

In triangulo rectangulo ex arcubus circulorum majorum
in superficie sphæræ, sinus unius acutorum angulorum
ad sinum totum est sicut sinus secundus reliqui acuti
ad sinum secundum lateris eum subtendentis.

Repetita descriptione 10 huius Demonstrandum est quod sinus
anguli $abc$ ad sinum totum est sicut sinus arcus $fe$ ad sinum
arcus $fb$. hoc modo per 6. præcedentis libri, et conversam
proportionem, sinus anguli $abc$ ad sinum totum, est sicut sinus
arcus $ac$ ad sinum arcus $ab$. Et per 7 eiusdem, et conversam
proportionem sinus arcus $ac$ ad sinum arcus $ab$ sicut sinus
arcus $cf$ ad sinum arcus $fb$. Quare per 11. 5. Euclidis sinus
arcus $cf$ ad sinum arcus $fb$ sic sinus anguli $abc$, ad sinum
totum quod erat demonstrandum.
\note{« $fe$ » à la troisième ligne et « $cf$ » plus bas : la main forme $c$ et $e$ presque de même ; les lettres sont données telles qu'elles se lisent, sans les accorder. Le second « $ac$ » est empâté d'encre.}
\note{figures dans la marge : un triangle lettré $a$, $b$, $c$, $d$, $e$, $f$ (les arcs $fbc$ et $abe$ se coupant en $b$) ; au-dessous une petite figure lettrée $l$, $a$, $b$, $e$, et à côté deux lettres encadrées, « $g$ » et « $h$ ».}

\subsection*{Propositio XXVIII.}

Item sinus secundus unius acutorum angulorum, ad
sinum totum est sicut quod sub sinu secundo arcus acutum
subtendentis, et sub sinu reliqui anguli continetur,
rectangulum ad quod ex sinu toto quadratum.

In descriptione præmissa, dico quod sinus secundus anguli $bac$
hoc est sinus arcus $fe$ ad sinum totum est sicut rectangulum
quod ex sinibus anguli $abc$ et arcus $fb$. qui est sinus secundus
arcus $bc$ ad quadratum sinus totius. Erit enim sinus anguli
$abc$ ad sinum totum per præcedentem sicut sinus arcus $ef$
ad sinum arcus $fb$, per 15. 6 Igitur rectangulum quod ex sinibus
anguli $abc$ et arcus $fb$ æquale est rectangulo quod ex sinu
toto sinuque arcus $ef$. Sed per primum 6. sicut sinus
arcus $fe$, ad sinum totum, sic quadrangulum ex sinu
toto sinuque arcus $fe$ ad quadratum sinus totius. Igitur
erit sinus arcus $fe$ ad sinum totum sic quadrangulum quod
\note{dans la marge, un cercle lettré $a$, $b$, $k$, $o$, $d$, $p$, $c$, $l$, $h$, $g$, $m$, $e$, $n$, $f$, avec la légende « Refertur ad propositionẽ Sequentem ».}

\note{page de droite de la vue 42, numérotée à l'encre « 41 » en haut à droite. Elle ne continue pas la page de gauche : celle-ci s'arrête sur « quadrangulum quod », au milieu de la démonstration de la proposition XXVIII, et celle-ci s'ouvre au milieu d'une application à l'astronomie. Les numéros à l'encre, 40 puis 41, se suivent pourtant. Le texte des pages numérotées 41 à 43 (vue 42, droite, à vue 45, gauche) se lit à la suite de la vue 47, gauche (« Si autem sumatur arcus $fbc$ pro horizonte recto » / « erit arcus $ac$ ascensio recta »), et la vue 47, droite, en reprend la fin (voir les notes des vues 45 et 47) : c'est une reconstitution de cette transcription, qui garde l'ordre des feuillets.}
erit arcus $ac$ ascensio recta, respondens arcui Zodiaci $ab$.
Atque quoniam per 17 præcedentis libri sinus secundus
arcus $bc$ ad sinum secundum arcus $ab$ est sicut sinus
totus ad sinum arcus $ac$. Hoc est sinus arcus $fb$ ad sinum
arcus $be$, sicut sinus totus ad sinum arcus $cd$, atque ex his
tria nota sint, scilicet sinus arcus $fb$, sinus arcus $be$ sinusque
totus, Iam et reliquum scilicet sinus arcus $cd$ noscetur.
Unde et arcus $cd$ dabitur. qui ablatus a quadrante relinquit
arcum $ac$ ascensionem scilicet rectam cognitam. Similiter et cæteræ
quæstiones absolventur, quæ faciles erunt horum libellorum
Theoremata intelligentibus. Lubet tamen hic demonstrare praxim
qua Iohannes Regiomontanus utitur in calculanda declinatione
et ascensione rectâ, per tabulas generales, nec non in habenda
differentia ascensionali. Non enim apertum est omnibus ingeniis
unde pendeant illius calculi regulæ. Atque pro tabula
declinationum sit Zodiacus $ab$ Æquator $ac$, Locus longitudinis
astri in Zodiaco punctum $b$. latitudo astri $bg$. Ita ut punctum $g$
sit verus astri locus. Declinatio arcus $gh$. Producatur arcus $gb$
usque ad æquatorem ad punctum $c$. Iam intranti in tabulam
declinationum generalem, cum longitudine astri occurrit primum
arcus $bc$. qui arcus est circuli latitudinis astri inter Zodiacum
et æquatorem comprehensus. Ex quo quidem arcu et latitudine
astri $bg$. sive adgregando sive subtrahendo, notus veniet arcus
$gc$ arcus \ill{} circuli latitudinis inter centrum astri et æquatorem\note{un mot surchargé et empâté après « arcus » ; on distingue « vi » au début, peut-être une abréviation, non résolue.}
Deinde occurrit in tabula numerus multiplicandus qui numerus
est sinus rectus anguli $gch$ posito sinu toto partium 100.000.
Unde \struck{primum} apud punctum tropicum quoniam talis angulus rectus
est, ponitur numerus ipse sinus totius, quandoquidem
anguli recti sinus est sinus totus, Per 6 autem præcedentis
libri, sinus totus ad sinum anguli $gch$ est sicut sinus arcus $cg$
ad sinum arcus $gh$. Quorum cum tria nota sint, scilicet sinus
totus 100.000, sinus anguli $gch$ qui est numerus multiplicandus
tabulæ, ac sinus arcus $cg$ Iam per regulam quatuor magnitudinum
noscetur et reliquum scilicet sinus arcus $gh$. Recte igitur
præcipit Iohannes in problemate dum iubet multiplicari sinum
arcus $gc$, in numerum multiplicandum in tabula repertum
hoc est multiplicari secundum in tertium et productum dividi
per primum, hoc est per sinum totum 100.000. Talis enim
divisio fit, dum in numero ex multiplicatione producto abiiciuntur
quinque figuræ ad dextram, quot sunt cyphræ in divisore,
hoc est in sinu toto: quæ operatio notissima est Arithmeticis,
Additur autem reliquarum figurarum numero unitas si
\note{dans la marge droite, un triangle lettré $g$, $b$, $a$, $h$, $c$, avec un second « $h$ » sous la base.}

\page{43}
\note{double page ; page de gauche. Elle continue la phrase de la fin de la vue 42 (« unitas si »).}
figuræ abiectæ denotant plus quam 50000. Quoniam fractio
quæ ultra integras partes ex divisione exit minor est dimidio.
cum abiectæ figuræ minus quam 50000. hoc est minus quam
dimidium divisoris denotant, et tunc negligitur. Maior
autem est dimidio cum relictæ figuræ plus dimidio dicto denotant, Et
tunc \struck{negligitur} pro fractione integra unitas adducitur hoc est integra
pars. Talis ergo numerus ex divisione sic facta progrediens, erit
quartum quod quærebatur, hoc est sinus ipsius arcus $gh$ hoc est
declinationis quæsitæ. Unde ex tabula sinus recti declinatio ipsa
cognita venit. Nec interest utrum sinus arcuum $cg$ $gh$ capiantur
alterius sinus totius quam sinus anguli $gch$ qui capitur respectu sinus
totius partium 100.000 quandoquidem in quatuor magnitudinibus
proportionalibus, non necesse est omnes quatuor esse eiusdem
colligantiæ, sed primam cum secundâ et tertiam cum quartâ.

Pro tabula autem ascensionum rectarum, repeto descriptionem
ultimæ propositionis huius libri, In qua quidem æquator sit $acd$
Circulus latitudinis astri $abe$. Circulus declinationis eiusdem
$fbc$, Circulus autem $fed$ incedens per polos ipsorum $abe$ $acd$
determinet ad $ae$ quadrantes. Punctum $b$ centrum astri,
Arcus $bc$ declinatio astri, Arcus $ac$ differentia transitus astri per
coeli medium. Itaque intranti tabulam ascensionum generalem
cum longitudine astri, occurrit primum arcus qui inscribitur
radix ascensionum. qui est quidem arcus æquatoris ab initio
arietis secundum ordinem signorum, usque ad circulum
latitudinis astri $abe$ computatus. Deinde occurrit numerus
multiplicandus qui nihil aliud est quam umbra versa
arcus $fe$ multiplicata in senarium posito gnomone partium
100.000, Et quidem per 31. huius præcedentem, sinus totus ad
sinum arcus $ac$ est sicut quadratum gnomonis ad quadrangulum
quod fit ex umbris versis arcuum $fe$ $bc$. Iam si dictum quadrangulum,
ducatur in sinum totum et productum dividatur
in quadratum gnomonis, exibit ex divisione sinus arcus $ac$.
Sed cum sinus totus hic supponatur partium 60000 et numerus
multiplicandus in tabula inventus sit umbra versa arcus $fe$ ducta
in 6, Atque quadratum gnomonis sit \uncertain{10000000.000} in idem
redit operatio, si ducatur numerus multiplicandus in umbram
versam arcus $bc$ declinationis, et a producto abiiciantur sex
figuræ quanquam canon iubet abiici quinque figuras ob id
videlicet quod numero multiplicando ablata sit una figura
ad alleviandum laborem. Et si abiectæ denotent plusquam
est 50000 adiicienda relicto numero unitas, ob causam sæpius
memoratam. Hoc pacto relictus est numerus is qui ex divisione
proditurus erat, hoc est sinus arcus $ac$. Unde et ipse arcus
$ac$ notus fiet. Cumque radix ascensionum per tabulam cognita
sit arcus æquatoris a puncto æquinoctiali verno usque ad punctum
\note{le nombre du carré du gnomon est écrit en chiffres serrés ; on y compte un 1 suivi de sept zéros, d'un point et de trois zéros, soit dix zéros, ce que demande le carré de 100.000 ; le compte des zéros reste douteux. « sæpius » est surchargé. Dans la marge gauche, en regard de « Pro tabula autem », un grand signe en croissant, « ( », qui n'est pas un mot.}
\note{figures dans la marge : en haut, un triangle lettré $f$, $b$, $e$, $c$, $a$ ; plus bas, la même figure lettrée $f$, $e$, $b$, $d$, $a$, $c$.}

\note{page de droite de la vue 43, numérotée à l'encre « 42 » en haut à droite ; elle continue la page de gauche (« usque ad punctum » / « a computatus »).}
a computatus: Iam addito vel ablato arcu $ac$ nunc noto,
cognitus veniet arcus æquatoris ab eodem æquinoctii puncto
usque ad punctum $c$ computatus: hoc est ascensio astri recta quæsita
quæ a circulo declinationis $fbc$ determinatur, Verum ipse arcus
$ac$ dria\note{abréviation de « differentia », écrite « dria » avec un trait.} scilicet transitus astri per medium coeli potest absque
auxilio numeri multiplicandi, et absque umbrarum tabula haben\note{écrit « haben » ; le sens attend « haberi ».}
sic. In descriptione præmissa per tabulam declinationum notus
erat arcus $cg$ latitudinis astri inter æquatorem $ac$ et centrum astri
$g$ comprehensus. Notus item fuit arcus $gh$ declinatio scilicet astri.
Atque ideo sinus secundi horum arcuum cogniti fient. Verum per
17. præmissi libri sinus secundus arcus $gh$ ad sinum secundum
arcus $gc$ est sicut sinus totus ad sinum secundum arcus $ch$
quandoquidem rectus est angulus $ghc$. Igitur cum ex his tria
nota sint, noscetur et reliquum scilicet sinus secundus arcus
$ch$. Et proinde ipse arcus $ch$ qui quærebatur notus erit, differentia
scilicet transitus astri per coeli medium, Unde ut prius per
radicem adscensionum, data fiet ascensio recta quæsita. Unde
pro ascensione recta satis erat radix ascensionum sola absque
numero multiplicando, et absque usu tabulæ foecundæ, quæ
tabula est umbrarum. Iohannes autem vir perspicacissimus
ideo adiecit numerum multiplicandum usumque tabulæ foecundæ,
\struck{quæ tabula est umbrarum}. Ut caveret divisionem quæ
necessaria est ei, qui per solam radicem ascensionum operatur,
quæ divisio cum scilicet divisor numerus diversas habet figuras
fastidiosior est numeranti, sic invenit modum, qui sola multiplicatione
perficeret, qui calculatoribus notior esse consuevit.

\note{avant « Pro regula », dans la ligne, le même signe en croissant qu'à la page de gauche.}
Pro regula tandem inveniendi ascensionalem differentiam, Sit
rursum in descriptione 31. præcedentis, æquator $acd$. Horizon
rectus $fbc$, obliquus $abe$, Meridianus $fed$. Eritque latitudo
regionis arcus $fe$. Centrum astri punctum $b$. Declinatio astri
arcus $bc$. Differentia ascensionalis eius arcus $ca$. Et quoniam
per 31. præced. quadratum gnomonis ad quadrangulum
factum ex umbris versis arcuum $fe$ $bc$: hoc est latitudinis loci,
et declinationis astri, Est sicut sinus totus ad sinum arcus
$ac$. Atque ex his quatuor tria nota sint scilicet quadratum
gnomonis, quadrangulum ex umbris factum, et sinus totus.
Ideo reliquum scilicet sinus arcus $ac$ notus veniet. Rectè
Igitur præcipit Iohannes, multiplicari umbras versas latitudinis
loci, ac declinationis astri alteram in alteram, \uncertain{ideoque}
quadrangulum produci prædictum. Tale enim quadrangulum
in sinum totum quem ipse supponit partium 60000 multiplicatum
ac in quadratum gnomonis. 10000.000.000 divisum exhibet sinum
arcus $ac$. Id autem fit dum prædictum quadrangulum in
\note{le nombre du carré du gnomon est écrit en groupes séparés par des points ; « 10000.000.000 », dix zéros ; le deuxième chiffre a la forme d'un « a », qui est le zéro de cette main (même forme dans « 100000 » à la vue 44). « quadratum » est surchargé.}
\note{dans la marge droite, deux figures : un triangle lettré $g$, $b$, $c$, $a$, $h$ ; plus bas, une figure lettrée $f$, $e$, $b$, $d$, $a$, $c$.}

\page{44}
\note{double page ; page de gauche. Elle continue la page de droite de la vue 43 (« quadrangulum in » / « senarium ducitur »).}
senarium ducitur: Et a producto 6 ad dextram figuræ abiiciantur,
reliquæ unitate adiectâ si abiectæ plusquam 50000 denotaverint.
Nam talis relictus numerus erit sinus arcus $ac$. Quare notescet
arcus $ac$, quæ differentia ascensionum quæsita est, Sic iam
demonstrata est praxis dictarum tabularum. ac simul declaratum est
qua ratione constructæ sint.

\subsection*{Demonstratio tabulæ beneficæ.}

Ad imitationem tabulæ foecundæ Iohannis de Monteregio
fecimus aliam tabulam quam beneficam appellavimus, quia eius
beneficio calculis quibusdam facilitas procuretur. Ut autem utriusque
tabulæ subiectum et speculatio patefiat. Esto triangulum rectilineum
$abc$ cuius angulus $abc$ rectus recta autem $bd$ perpendicularis ad
basim $ac$. Ponaturque linea $ab$ sinus totus et gnomo proiiciens
umbram $bc$ qui sinus et gnomon supponitur partium 100.000.
Eritque angulus $a$ cum quo intramus tam in tabulam foecundam
quam in tabulam beneficam, et etiam in tabulam sinus recti, Per
quem ingressum ex tabula foecunda habetur linea $bc$ quæ est
umbra versa ex tabula benefica habetur linea $ca$ radius scilicet
contingens apicem gnomonis scilicet punctum $a$ cum $c$ puncto
extremo umbræ, ex tabula sinus habetur linea $bd$ sinus scilicet
anguli $a$ cuius sinus secundus est linea $da$ sinus scilicet anguli
$abd$ quod est complementum anguli $a$ quandoquidem conjuncti
faciunt angulum rectum. Quoniam igitur triangula $abd$ $bcd$
partialia sunt per 8. 6 Eucl. similia inter se et totali triangulo
$acb$ propterea erunt lineæ $da$ $ab$ $ca$ continue proportionales.
Quare si quadratum ipsius $ab$ dividatur per lineam $ad$ numerus
quotiens erit linea $ca$. Item erit sicut $ab$ linea ad ipsam
$bd$ sic $ac$ ad ipsam $cb$. Quare si ducatur $ac$ in $bd$ et productum
secetur in $ab$, quod fiet abiectis de producto quinque ad
dextram figuris quot scilicet sunt zyphræ in numero 100000
numerus quotiens erit $bc$ umbra versa. Contra erit sicut $bd$
ad ipsam $ab$, sic et $bc$ ad ipsam $ca$. Quare si ducatur $ab$
in $bc$ quod fiet applicando ad dextram numeri $bc$ quinque
zyphras. Et productum secetur in $bd$, numerus quotiens erit
linea $ca$. Item si ponatur angulus $a$ graduum 30 tunc $bd$
sinus erit dimidium ipsius $ab$ sinus totius. Sed sicut $ab$ ad
ipsam $bd$ sic $ac$ ad ipsam $cb$. Ergo tunc $bc$ umbra erit dimidia
numeri $ca$ tabulæ beneficæ: Et quadratum ipsius $ac$ quadruplum
ad quadratum ipsius $bc$, Et ideo quadratum ipsius $ac$ sesquitertium
ad quadratum ipsius $ab$. Unde tunc si ponatur $ab$
semidiameter sphæræ, lineam iam $ac$ erit latus cubi circumscripti
\note{dans la marge, un triangle rectangle lettré $a$, $b$, $c$, avec la perpendiculaire $bd$ sur $ac$. Plusieurs mots de la page sont soulignés (« similia », « angulus », « tunc », « erit », « ab ») sans raison apparente ; le soulignement n'est pas reproduit.}

\note{page de droite de la vue 44, numérotée à l'encre « 43 » en haut à droite ; elle continue la page de gauche (« cubi circumscripti » / « a sphæra »).}
a sphæra: sed hoc nihil ad rem. Adhuc si ponatur angulus
$a$ graduum 45. hoc est dimidium recti $ad$ $db$ $dc$ erunt inter
se æquales. Unde $bc$ umbra erit æqualis gnomoni suo $ab$ et
$ac$ numerus tabulæ beneficus duplus ipsius $ad$ sive $db$ qui
sunt sinus primus et secundus anguli $a$. Et tunc si ponatur $ab$
semidiameter circuli erit $ac$ latus quadrati in circulo descripti
sed nihil ad rem. Demum si ponatur angulus $a$ graduum 60
tunc angulus $c$ fiet graduum 30, et tantundem angulus $abd$.
Quare $ad$ erit dimidium ipsius $ab$. sed sicut $ba$ ad ipsum $ad$ sic
est $ca$ ad ipsum $ab$. Ergo $ab$ tunc dimidium ipsius $ac$. Unde
cum $ab$ ponatur partium 100000 erit eo casu $ac$ numerus
tabulæ beneficæ partium 200.000 Postremo si ponatur angulus
$a$ graduum 90 hoc est rectus, tunc in tabula sinuum ipse sinus
$bd$ est iam ipse sinus totus $ab$. Nam in eo casu linea $bd$ coünitur ipsi
$ab$. Tunc vero tam in tabula Foecunda quam in tabula benefica
ibi umbra versa et hic radius in infinitum abeunt, quoniam scilicet
$bc$ umbra et $ac$ radius sunt æquidistantes, et necubi quamvis in
infinitum producti concurrunt. Ex quibus satis constat deffinitio et
fabrica ipsarum tabularum Foecundæ et Beneficæ. Nam ex tabula sinuum
per regulas proprias compositæ eliciuntur ambæ.

\note{ici, dans la ligne, le signe en croissant déjà vu aux vues 43 ; « quoad » est écrit « qual » corrigé par « oa » en interligne.}
Porro quoad usum tabulæ
beneficæ, nam foecundæ calculus fuit dudum demonstratus: hoc accipe.
Sint in superficie sphæræ duo quadrantes circulorum, majorum $abe$
$acd$ inter se inclinatorum, Et a puncto $f$ polo scilicet ipsius $ae$ descendant
duo quadrantes $fed$, $fbc$ Eruntque anguli apud $c$ $d$ $e$ puncta
recti, Sintque ex triangulo sphærali $abc$ cogniti arcus $ab$ $bc$, quæro
ex his arcum $ac$. Sic cum arcu $bc$ intro in tabulam Beneficam, Et
excipio numerum multiplicandum qui sit $g$. Et quoniam sinus
secundus arcus $bc$ est ipse sinus arcus $fb$ Erunt iam sicut dudum
ostensum est sinus $fb$ arcus sinus totus, atque $g$ continue proportionales,
Sed per demonstrata Menelai sicut sinus ipsius $fb$ ad sinum
totum arcus $fc$ sic sinus ipsius $be$ ad sinum ipsius $cd$. Igitur
sicut sinus ipsius $be$ ad sinum ipsius $cd$, sic erit sinus totus ad
ipsum $g$. Quamobrem multiplico ipsum $g$ in sinum ipsius $be$
Et productum divido in sinum totum quod fit abiiciendo quinque
\struck{zyphras} figuras\note{« figuras » est écrit au-dessus de « zyphras » biffé.} de producto ad dextram quot sunt zyphræ in sinu toto.
Et habebo in quotiente sinum arcus $cd$, quo ablato a quadrante
superest mihi arcus $ac$ quæsitus. Atque hoc quidem pacto ex $ab$ arcu
Eclipticæ a sectione æquinoctii incepto Et ex $bc$ declinatione cognoscam
$ac$ rectam ascensionem

\note{même signe en croissant avant « Item ».}
Item mutatis terminis ex $ab$ arcu
circuli latitudinis inter stellam $b$ et æquatorem $acd$ clauso Et ex
$bc$ declinatione stellæ datis: discam arcum $ac$ æquatoris qui dicitur
\note{« apud $c$ $d$ $e$ puncta » : écrit « apud cde puncta ». Dans la marge droite, deux triangles rectangles lettrés $a$, $b$, $c$, $d$ (la perpendiculaire $bd$) ; plus bas, un gros point et un rectangle portant $a$ en haut et $b$ en bas.}

\page{45}
\note{double page ; page de gauche. Elle continue la page de droite de la vue 44 (« qui dicitur » / « differentia »).}
differentia stellæ transitus per coeli medium, et qui ad sciendam
eius ascensionem rectam usu venit

\note{signe en croissant avant « Adhuc », et de même avant « Si autem », « Quod si » et « Denique » plus bas.}
Adhuc variando iuxta necessitatem
propositi terminos ex latitudine ortus $ab$ et ex declinatione
$bc$ stellæ propositæ in puncto $b$ positæ datis eliciam arcum $ac$ diff\textsuperscript{ia}
ascensionalis eiusdem stellæ, Et sic deinceps in cæteris quæstionib\textsuperscript{9}
quæ ad hanc arcuum descriptionem redigi possunt.

Si autem
in eodem triangulo sphærali $abc$ datus sit angulus $a$ hoc est
arcus $de$ et arcus $bc$. tunc ex his scitur arcus $ab$ sic. Cum
arcu $fe$ intro tabulam beneficam, et numerus multiplicandus
ibi inventus sit $g$. Et quoniam sinus secundus arcus $fe$ est sinus
arcus $ed$, Erunt sicut prius sinus $ed$ arcus, sinus totus et $g$
continuæ proportionales. Sed sicut ostendit idem Menelaus
sicut est sinus ipsius $ed$ ad sinum totum sic iam sinus ipsius $bc$
ad sinum ipsius $ab$. Ergo et sicut sinus ipsius $bc$ ad sinum
$ab$, sic sinus totus ad ipsum $g$. Quamobrem multiplico $g$ in
sinum ipsius $bc$ et productum partior in sinum totum, quod
fit abiectis quinque dexterioribus figuris et exibit ex divisione
sinus ipsius $ab$ arcus quæsiti. Tali calculo ex arcu $bc$ declinationis
et ex maxima declinatione $ed$ datis comperitur arcus $ab$
Eclipticæ, qui arcui ascensionis $ac$ respondet. Et similiter
faciam in aliqua simili quæstione ut si ex angulo $a$ complemento
latitudinis regionis et ex arcu $bc$ declinationis stellæ,
Si ex angulo $a$ complemento latitudinis regionis et ex arcu
$bc$ declinationis stellæ, in puncto $b$ positæ velim noscere
arcum $ab$ horizontis inter æquatorem et stellam exorientem
qui latitudo ortus vocatur. Sicut dictum est.\note{la même proposition, « si ex angulo $a$ complemento latitudinis regionis et ex arcu $bc$ declinationis stellæ », est écrite deux fois de suite, sans biffure.}

Quod si dati
sint duo arcus $ac$ $cb$ tunc per eos inveniam arcum $ab$. Dabuntur
enim arcus $cd$ $fb$ quæ sunt datorum complementa \struck{Et tunc per
eos inveniam arcum $ab$}. Et ex his per regulam declinationis
dabitur arcus $be$, Et eius complementum $ab$ quæsitum.

Denique
si ex arcubus $ab$ $bc$ datis quæratur arcus $cd$. Tunc intrabis tabulam
beneficam cum arcu $be$, et numerus inventus sit $g$. Et quia
sinus secundus ipsius $be$ est sinus ipsius $ab$. Idcirco erunt ut
prius, sinus ipsius $ab$, sinus totus, et ipsum $g$ continuæ
proportionales. Verum ut ostendit Menelaus sicut sinus
ipsius $ab$ ad sinum totum, sic iam sinus ipsius $bc$ ad sinum
ipsius $cd$. Igitur erit sicut sinus ipsius $bc$ ad sinum
ipsius $cd$. sic sinus totus ad ipsum $g$. Quare multiplico
$g$ in sinum ipsius $bc$ et productum partior in sinum totum,
quod fit abiectis quinque figuris, vel quot zyphræ sunt in numero
\note{la page s'arrête sur « in numero » ; la suite ne se lit pas à la vue 45, droite, mais en tête de la vue 47, droite (« sinus totius, Et habeo in quotiente sinum arcus $cd$ quæsiti »).}
\note{« prius » : les lettres du milieu sont surchargées, peut-être sur « primus ». Dans la marge, deux fois la même figure : un triangle lettré $f$, $e$, $b$, $a$, $c$, $d$, et au-dessous un trait horizontal marqué $g$.}

\note{page de droite de la vue 45, numérotée à l'encre « 44 » en haut à droite. Elle ne continue pas la page de gauche : elle reprend, en les répétant, les deux dernières lignes de la page de gauche de la vue 42 (« toto sinuque arcus $fe$ ad quadratum sinus totius. Igitur erit sinus arcus $fe$ ad sinum totum sic quadrangulum quod ») et achève la proposition XXVIII. Les pages numérotées 41 à 43 (vues 42, droite, à 45, gauche) forment un texte suivi sur les tables de Regiomontanus et la « tabula benefica », placé au milieu de la proposition XXVIII.}
ex sinu toto sinuque arcus $fe$ ad quadratum sinus totius. Igitur
erit sinus arcus $fe$ ad sinum totum sic quadrangulum quod
ex sinibus anguli $abc$ et arcus $fb$ ad quadratum sinus totius.
Quod est propositum.

\subsection*{Propositio XXVIIII.}

In omni triangulo ex arcubus circulorum majorum
in superficie sphæræ constituto, sinus versus anguli cuiuslibet
ad differentiam duorum sinuum versorum quorum unus
est lateris eum angulum subtendentis, alter vero differentiæ
duorum arcuum ipsum angulum continentium, Est sicut
quadratum sinus totius ad id quod sub sinibus arcuum
eorundem continetur rectangulum.

Sit in superficie sphæræ ex arcubus circulorum majorum
constructum utcunque triangulum $abc$. Completoque circulo
$adb$ sub polis $ab$ circuli describantur, majores quidem \uncertain{2}
$def$ $geh$, minores vero $kcn$, $ocp$, Item quadrantes circulorum
majorum a polis ad circulos majores $acl$, $bcm$. Rursus in alia
descriptione (confusionis tollendæ gratia) describantur horum circulorum
diametri axes, sinusque. Ac primum descripto utcunque circulo
$abd$ describantur aliorum diametri, quæ sunt communes eorum
cum $abd$ circuli plano sectiones hoc pacto. Punctum $x$. Sit centrum
sphæræ et circulorum majorum, Recta $dxf$ diameter circuli $def$,
Recta $gxh$ diameter circuli $geh$. Recta $kyn$ diameter circuli $kcn$
Recta $orp$ diameter circuli $ocp$ sese invicem secantes in puncto \uncertain{$r$}. Item
rectæ $xb$ $xa$ axes ipsorum $geh$ $def$ circulorum sive semidiametri
circulorum $bcm$ $acl$, Adhuc lineæ $au$ $kq$ perpendiculares ad $xb$,
Linea $kt$ perpendicularis ad \uncertain{$os$}. Quibus peractis patet quod linea
$au$ est sinus rectus arcus $ab$. Linea $ky$ est sinus rectus arcus $ak$.
et etiam arcus $ac$. quandoquidem æquales quoniam a polo ad
periphæriam circuli $ocp$. Linea $kq$ sinus arcus $bk$ qui est differentia
arcuum $ak$ $ab$. Unde et linea $br$ erit sinus versus arcus $bo$ sive
$bc$. linea $bq$ sinus versus arcus $bk$ differentiæ prædictæ. Linea $qr$
differentia talium sinuum versorum. linea $ks$ sinus versus arcus $kc$. Quæ
omnia patent per primam præcedentis libelli, Ponatur itaque recta $dz$
sinus versus arcus $db$. sive anguli $bac$. Constabitque demonstrandum
esse quod $dz$ sinus versus anguli $bac$ ad $ry$ driam\note{abréviation de « differentiam », écrite « driam » avec un trait.} sinuum
versorum, arcubus $bc$ $bk$ debitorum est sicut quadratum sinus totius.
quod fit ex $dx$ in $xa$, ad rectangulum quod fit $au$ $ky$. qui
\note{« majores quidem 2 » : un signe en forme de 2 après « quidem ». Dans la marge droite, quatre figures : deux sphères traversées d'arcs, sans lettres ; une troisième lettrée $a$, $b$, $k$, $o$, $d$, $g$, $l$, $m$, $e$, $c$, $f$, $n$, $h$, $p$ ; une quatrième, en coupe, lettrée $a$, $b$, $u$, $k$, $o$, $d$, $q$, $r$, $s$, $t$, $p$, $y$, $z$, $h$, $x$, $g$, $n$, $f$. Le texte s'arrête sur « qui » au bas de la page.}

\page{46}
\note{double page ; page de gauche. Elle continue la page de droite de la vue 45 (« au ky. qui » / « scilicet sunt »).}
scilicet sunt sinus arcuum, $ab$ $ac$ recti. hoc modo. Quoniam
per 14. secundi sphæricorum elementorum arcus $dl$ $ks$ sunt
similes, Et in similibus arcubus sinus sunt semidiametris
proportionales, Et ideo sicut linea $dz$ ad lineam $ks$ sic linea
$dx$ ad lineam $ky$. Item quoniam æquiangulum est triangulum
$kts$. triangulo $\supset rs$. Et triangulum $\supset ts$, triangulo $\supset yx$. Et
triangulum $\supset yx$ triangulo $aux$ ideo triangula $kts$, $aux$
æquiangula sunt. Quare per 4. 6. Euclidis, linea $xa$ ad lineam
$au$ sicut linea $sk$ ad $kt$. Verum proposita linea $ks$ media ratio
$dz$ ad $kt$ vel $qz$ componitur ex rationibus $dz$ ad $ks$
et $ks$ ad $kt$ Et ideo ex rationibus $dx$ ad $ky$. Et $xa$ ad $au$
quæ sunt illis eædem, Et per 24. 6. Elem. Euclidis ratio quadrati
sinus totius quod fit ex $dx$ $xa$ ad quadrangulum quod fit
ex $au$ $ky$ componitur ex rationibus $dx$ ad $ky$, et $xa$ ad $au$
iisdem, Igitur ex æqua proportione Erit linea $dz$ ad \uncertain{$aqz$}
sicut quadratum sinus totius ad rectangulum $au$ $ky$. quod quidem
proponebatur demonstrandum.
\note{« $\supset$ » : une lettre de figure en forme de c retourné, notée ici par le signe le plus proche, distincte de $c$ et de $z$ ; elle figure aussi, près de $r$, sur la figure de la vue 45. Le premier « $\supset rs$ » est surchargé. Les lettres « $aqz$ » sont repassées à l'encre et ne nomment aucune ligne de la figure ; l'énoncé de la vue 45 demande ici $ry$. « $qz$ » (après « vel ») est lu tel quel.}
\note{dans la marge, trois figures qui reprennent celles de la vue 45 : une sphère traversée d'arcs, sans lettres ; une sphère lettrée $b$, $a$, $k$, $p$, $o$, $d$, $c$, $h$, $g$, $n$, $e$, $m$, $f$ ; une coupe lettrée $a$, $b$, $u$, $k$, $o$, $d$, $p$, $q$, $s$, $t$, $y$, $\supset$, $h$, $x$, $g$, $n$, $f$.}

\subsection*{Propositio XXX.}

Sinus arcus alicuius ad sinum eius secundum est sicut
gnomo ad umbram rectam eiusdem arcus, sicut autem
secundus ad sinum dicti arcus, sic est gnomo ad
umbram versam eiusdem arcus.

In quadrante circuli $ab$ capiatur quantusvis arcus $bc$, Eiusque
complementum $ca$. Ductisque semidiametris $ad$ $db$, et $ce$ perpendiculari
ad $db$. Erit per primam præcedentis libri recta \uncertain{$ca$}
sinus rectus arcus $bc$. Recta vero $ed$ sinus secundus eiusdem arcus
$bc$ coniungatur $cd$ et producatur ad $g$ punctum quodvis. Producatur
item $ad$ $bd$ semidiametri. Compleaturque rectangulum
$dfgh$. Itaque posito gnomone $df$ umbra recta respondens arcui
$bc$ erit recta $fg$. Posito autem gnomone $dh$ umbra versa
eiusdem arcus erit $gh$. Quibus declaratis demonstrandum est quod
$ce$ ad $ed$ est sicut $df$ ad $fg$. Item quod $de$ ad $ec$ est sicut
$dh$ ad $hg$. Constat autem utrumque ex similitudine
triangulorum $ced$ $dfg$, $ghd$. et proportione laterum.
\note{figure dans la marge : un quadrant lettré $a$, $c$, $b$, $e$, $d$, $h$, $f$, $g$, avec le rectangle $dfgh$.}

\subsection*{Corollaria.}

Unde umbra recta cuiusvis arcus est et umbra versa complementi
eiusdem arcus. Item Dimidio quadrantis debita umbra tam recta
quam versa æqualis est suo gnomoni. Item gnomon semper est
medius proportionalis inter umbras duorum arcuum quadrantem
integrantium, seu rectas seu versas.

\subsection*{Propositio XXXI.}

In triangulo rectangulo ex arcubus circulorum majorum in
superficie sphæræ, constituto, sinus totus ad sinum arcus unum

\note{page de droite de la vue 46, numérotée à l'encre « 45 » en haut à droite ; elle continue la page de gauche (« ad sinum arcus unum » / « ex acutis »).}
ex acutis subtendentis est sicut quadratum gnomonis
ad rectangulum comprehensum sub umbris versis quorum
una complemento reliqui acuti, altera vero eundem
acutum subtendenti debetur.

Repetita descriptione 27. huius in triangulo $abc$ demonstrandum
est quod sinus totus ad sinum arcus $ca$ subtendentis angulum
acutum $abc$ est sicut quadratum gnomonis ad rectangulum
quod fit ex duabus umbris versis, quarum una debetur arcui
$fe$ quod complementum est anguli $bac$ Altera \uncertain{respondet} arcui
$bc$ subtendenti reliquum acutum angulum $bac$. Quod sic
ostenditur. Per præcedentem sicut gnomo ad umbram versam
arcus $fe$ sic sinus arcus $de$ ad sinum arcus $ef$. Et rursum
per eandem sicut gnomo ad umbram versam arcus $bc$ sic
sinus arcus $fb$ ad sinum arcus $bc$. Sed per 24. 6 Elem. Eucl.
Ratio quadrati gnomonis ad quadrangulum ex umbris versis arcuum
$fe$ $bc$ componitur ex rationibus duabus videlicet gnomonis ad umbram
versam arcus $fe$. Et gnomonis ad umbram versam arcus $bc$. Quare
Et eadem ratio componetur ex rationibus duabus quæ sunt prædictis
æquales, scilicet sinus arcus $de$ ad sinum arcus $ef$ Et sinus arcus
$fb$ ad sinum arcus $bc$. Per 4. vel 5. huius. Ex iisdem rationibus
duabus scilicet sinus arcus $de$ ad sinum arcus $ef$, et sinus arcus
$fb$ ad sinum arcus $bc$ componitur ratio sinus arcus $da$ ad sinum
arcus $ac$. Igitur ex æqua proportione erit sinus arcus $da$ hoc
est sinus totus ad sinum arcus $ac$. Igitur ex æqua proportione
erit sinus arcus $da$ ad sinum arcus $ac$. Iam sicut quadratum
gnomonis ad rectangulum contentum sub umbris versis arcuum
$fe$ $bc$, et id proponebatur demonstrandum.
\note{« Altera » est suivi d'un mot abrégé, lu « rndet » avec un trait, offert comme « respondet ». La phrase « Igitur ex æqua proportione erit sinus arcus $da$ … ad sinum arcus $ac$ » est écrite deux fois de suite ; à la seconde, « hoc est sinus totus » est ajouté en interligne au-dessus de « $da$ ». Dans la marge, un triangle lettré $f$, $e$, $b$, $a$, $d$, $c$.}

\subsection*{Corollaria.}

Hinc manifestum est, quod si in duobus triangulis sphæralibus
orthogoniis, duo latera quæ circum rectos angulos æqualia fuerint
tunc rectangula quæ ex umbris versis reliquorum laterum, quæ
circum rectos in umbras versas debitas complementis subtensorum
angulorum producuntur, erunt et invicem æqualia. Et contrario.

\subsection*{Propositio XXXII.}

In triangulo rectangulo ex arcubus magnis in superficie sphæræ,
constituto, sinus totus ad sinum arcus unum ex acutis angulis
subtendentis est sicut umbra versa reliquo acuto angulo debita
ad umbram versam lateris dictum acutum respicientis.

Hoc est in eadem descriptione \ill{} quod sinus totus scilicet
arcus $da$ ad sinum arcus $ac$ est sicut umbra versa arcus $de$
ad umbram versam arcus $bc$. Sic per ultimum corollarium 30
præcedentis ex umbra versa arcus $fe$ sui complementi producit
\note{après « descriptione », un mot abrégé d'environ cinq lettres (« ontaz » ou approchant), non lu. Dans la marge, la même figure lettrée $f$, $e$, $b$, $a$, $d$, $c$.}

\page{47}
\note{double page ; page de gauche. Elle continue la page de droite de la vue 46 (« sui complementi producit » / « quadratum gnomonis »).}
quadratum gnomonis. Sed per præcedentem quadratum gnomonis
ad rectangulum umbrarum versarum arcuum $fe$ $bc$ fuit sicut
sinus arcus $da$ ad sinum arcus $ac$. Igitur et sicut quod fit
ex umbra versa arcus $de$ in umbram versam arcus $fe$, ad
rectangulum quod fit ex umbris versis arcuum $fe$ $bc$ sic sinus arcus
$da$ scilicet sinus totus ad sinum arcus $ac$. Cumque per primam
6. Euclidis, rectangulum umbrarum versarum arcuum $de$ $fe$,
ad rectangulum umbrarum versarum arcuum $fe$ $bc$. sit sicut
umbra versa arcus $de$ ad umbram versam arcus $bc$. Erit iam
sicut sinus totus ad sinum arcus $ac$ sic umbra versa arcus
$de$ ad umbram versam arcus $bc$. Quod fuit demonstrandum.
Ex eadem omnia concludere potes de rectis umbris quæ
talium arcuum complementis debentur.
\note{dans la marge, un triangle lettré $f$, $e$, $b$, $d$, $c$, $a$ ; plus bas, le long de la reliure, deux « $f$ » isolés, peut-être des lettres d'une figure prise dans le pli.}

\subsection*{Scholium.}

Ex his igitur demonstrationibus absolvi possunt omnes quæstiones
quæ fieri consueverunt circa sphæralia triangula quæque in
astronomia pertinent ad arcus circulorum in primo coelo intellectorum,
hoc in concava superficie primi mobilis descriptorum
Sed nemo harum speculationum scientiam perfectam habens nesciet
theoriam ad praxim atque demonstrationem ad calculum reducere.
Quæ res ut magis pervia fiat lectori exempla nonnulla sunt adducenda
ut circa declinationes et adscensiones.

Et primum si in descriptione præcedentis circulus $acd$ sit æquator
circulus $abe$ sit Zodiacus. Circulus $fed$ sit colurus incedens per
utrosque eorum polos punctum $a$ punctum æquinoctii, Arcus $de$
maxima Zodiaci declinatio, Punctum $b$ locus astri in Zodiaco
Circulus $fbc$ circulus declinationis, arcus $bc$ declinatio astri,
ex loco autem astri cognito et maxima declinatione datâ hoc
ex arcubus $ab$ $de$ quæratur declinatio astri hoc est arcus $bc$.
absolvetur quæstio sic, cum per 6. \struck{huius} præmissi libri sicut est
sinus totus ad sinum arcus $de$ sic \uncertain{sit} sinus arcus $ab$ ad sinum
arcus $bc$. Atque ex his quatuor tria cognita sint noscetur et
reliquum per regulam quatuor magnitudinum, hoc est
cum sinus totus datus sit sinusque arcus $de$ dati datus, nec
non sinus arcus arcus $ab$ dati propter cognitum astri locum
datus Iam et reliquus sinus arcus $bc$ quæsiti dabitur, Et perinde
arcus $bc$ quæsitus dabitur. Unde si ducatur secundum in
tertium et productum dividatur per primum exibit ex divisione
quartum ex doctrina 14. 6. Elem. Eucl. Hoc est si ducatur sinus
arcus $de$ in sinum arcus $ab$ et productum dividatur in sinum
totum exibit in divisione sinus arcus $bc$ Unde et arcus $bc$ quantus
notus erit.

\note{signe en croissant avant « Si autem ».}
Si autem sumatur arcus $fbc$ pro horizonte recto
\note{la page s'arrête ici, au bas du feuillet, la phrase en suspens ; elle se poursuit, selon toute apparence, en tête de la vue 42, droite (« erit arcus $ac$ ascensio recta, respondens arcui Zodiaci $ab$ »). Dans la marge, en face du milieu de la page, le même triangle lettré $f$, $e$, $b$, $d$, $c$, $a$.}

\note{page de droite de la vue 47, numérotée à l'encre « 46 » en haut à droite. Elle ne continue pas la page de gauche, qui s'arrête sur « Si autem sumatur arcus $fbc$ pro horizonte recto » : celle-ci s'ouvre sur la fin d'un autre calcul (« sinus totius, Et habeo in quotiente sinum arcus $cd$ quæsiti »), celui de la plus grande déclinaison $de$ tirée de l'arc $ab$ et de la déclinaison $bc$.}
sinus totius, Et habeo in quotiente sinum arcus $cd$ quæsiti. Hoc
ergo pacto ex $ab$ arcu Eclipticæ et ex eius declinatione $bc$
nanciscemur maximam declinationem $de$. Ex quibus
quidem manifestum est quod omnia \ill{} quæ Iohannes
de Monte Regio per tabellam foecundam vitato divisionis
fastidio per multiplicationem elaborabat, hic per nram\note{abréviation de « nostram », écrite « nram » avec un trait.}
beneficam haud difficilius supputari possunt sed practica
huiusmodi regularum exempla inferius una cum tabellis
exponemus. Hæc in Arce Apollinari dum cum Iohanne
Vigintimillio Hieracensium Marchione degeremus olim
mense Augusto 1530 speculabamur. Nunc adnectemus
his quædam veterum opuscula ad eandem sphæricam
theoriam spectantia.

\subsection*{D. Octavio Spinulæ Quæstori ærario Siciliæ, Regioque Consiliario, Maurolycus. Salutem.}

Ex quo Sphæricorum nostrorum volumen Octavi generose
vir Messanæ excudi coeperat fluxit iam annus tertius
Et inter varia tum impedimenta bellorumque suspiciones
quæ nos turbabant, non dabatur ullus speculationibus
locus quæ secessum mentis et otium imprimis requirunt:
fecit amor erga me tuus, immo eximia tua in litteras
affectio, ut repetito labore opus hoc egregium absolveretur
Tantum igitur tibi sphæra nostra debet, ac debere se
fatetur, quantum Atlas Herculi coelum humeris sustinenti
debuit. Absolutis itaque Theodosij Menelai ac Maurolyci
tui Sphæricis, subiungemus ea quæ ad sphæram mobilem
motumque primum pertinent, quippe quæ sphæricorum
demonstratis innituntur, et prima sunt disciplinæ
syderalis rudimenta. Quod si te nunc Reip. cura et onus
quod sustines occupatum tenet, dabitur ut spero \uncertain{quandoque}
ocium, ut artibus ingenuis intendas animum et opera
nostra utaris. Vale.
\note{« omnia » est suivi d'un court mot biffé, non lu (marqué illisible). « quandoque » est écrit en abrégé (« qdoq » avec un signe final).}

\subsection*{Ad eundem Hexastichon.}

\begin{quote}
Disturbata diu varijs monumenta procellis
\end{quote}

\begin{quote}
Ad metam properant, te duce nostra suam.
\end{quote}

\begin{quote}
Haud aliter quondam ne machina tanta periret,
\end{quote}

\begin{quote}
Herculis auxilio sydera fulsit Atlas.
\end{quote}

\begin{quote}
Si quidquam Octavi poterit mea pagina, laudem
\end{quote}

\begin{quote}
Delebunt meriti, tempora nulla tui.
\end{quote}

Messanæ. Kal. Junij M.D.LVII.
\note{les vers vont à la ligne un par un, les pairs en retrait. « Disturbata » et « procellis » sont surchargés ; le « D » de la date est repassé.}

\page{48}
\note{double page ; page de gauche. Elle ne continue pas la page de droite de la vue 47, qui s'achève sur la date de la dédicace : un nouveau texte commence ici, en tête de page. Les numéros des propositions sont écrits dans la marge gauche, en regard de l'énoncé.}

\section*{Autolyci de sphæra quæ movetur ex traditione Maurolyci. Liber.}

Sphæræ puncta æqualiter ferri dicuntur quæcumque æquali
tempore æquales ac similes prætereunt periphærias. At
si in linea aliqua delatum aliquod punctum æqualiter
binas transierit lineas eandem habebit rationem tempus
tempus, quibus singulas transit lineas, quam linea
ad lineam.\note{« rationem tempus tempus » : « tempus » en fin de ligne puis « tempus » en début de ligne suivante, sans « ad ».}

Hoc est peracta spatia temporibus proportionalia sunt.

\textbf{Propositio 1.} Si æqualiter sphæra volvatur circa suum axem cuncta
quæ in superficie sunt sphæræ puncta, præter polos
circulos describent parallelos et ad axem rectos et
eosdem cum sphæra polos habentes.

Nam tales circuli describuntur per rectos a punctis ad
axem super quo sphæra versatur perpendiculares, et
ideo per 9 primi Theodosij habent dictum axem communem
et polos communes.

\textbf{Propositio 2.} Si sphæra volvatur æqualiter circa suum axem
cuncta quæ in superficie puncta sunt sphæræ similes
periphærias circulorum parallelorum in quibus
feruntur in tempore eodem præteribunt.

Nam si duo puncta sint in eodem parallelo. Constat
propositum per assumptam in principio petitionem.
Si in diversis parallelis autem, constabit propositum
adducta 15. 2. sphæricorum Theodosij.

\textbf{Proposito 3.} Si æqualiter sphæra evolvatur circa suum axem
quas in tempore eodem periferias transmittent
puncta quædam in circuli parallelis per quos
feruntur eæ ipsæ similes erunt.

Hæc est conversa præcedentis et converso modo demonstratur.

\textbf{Propositio 4.} Si in sphærâ major circulus manens separet
id quod apparet de sphæra ab eo quod non
apparet, sitque ad rectos angulos axi sphæræ
super quo movetur. Nullum punctum superficiei
sphæræ oritur, nullum occidit, sed quæ sunt
in hæmisphærio apparenti semper apparent
Quæ vero in latenti semper occultantur.

Nam talis circulus manens est communis limes talium
hæmisphæriorum, et per primam huius puncta singula
\note{« Proposito 3. » est écrit ainsi dans la marge. La fin de « hæmisphæriorum » est surchargée. La page s'arrête sur « singula » ; la page de droite de la même vue ne la continue pas.}

\note{page de droite de la vue 48, numérotée à l'encre « 47 » en haut à droite. Elle ne continue pas la page de gauche (le livre d'Autolycus, qui s'arrête sur « puncta singula ») : elle s'ouvre au milieu d'une suite de problèmes numérotés sur les tables, avec des exemples chiffrés.}
qui quidem arcus est horizontis inter æquinoctialem, et solem
in tali horizonte orientem.

\note{un « 7 » est écrit au-dessus de la fin de la première ligne de ce paragraphe ; au paragraphe suivant, un « 8. » dans la marge gauche : apparemment les numéros des problèmes.}
Data item loci latitudine atque stellæ cuiuspiam declinatione
differentiam ascensionalem tali stellæ debitam perscrutari.

Cum tali declinatione data sume numerum ex tabula benefica
quem multiplica in sinum secundum latitudinis ortivæ talis
stellæ ad talem locum: Et a producto abiice quinque dextrorsum
figuras: Et relicto sinui attribue arcum, namque is de quadrante
summotus residuabit ascensionem quæsitam. Exempli gratia
Ad dicti loci latitudinem gr. 38 ex tabula benefica excipio
numerum 109038 quem multiplico in 86251. sinum scilicet
secundum latitudinis ortivæ, talium declinationis et loci,
scilicet graduum 30. m. 24. Et produco 940.463.6538.
Hinc abiectis quinque ad dextram figuris, restat numerus
94046. cui sinui debetur arcus gr. 70. m. 8. cuius comp\textsuperscript{tum}
gr. 19. m. 52. est differentia ascensionalis quæsita, quæ
quidem pro declinatione Septemtrionis adiecta pro declinatione
ad austrum subtracta quadranti residuat arcum semidiurnum.
\note{dans « 940.463.6538 », les chiffres « 53 » sont surchargés.}

\textbf{8.} In horizonte proposito astri cuius datur altitudo ac
declinatio, distantia a meridiano vestigatur.

Ponatur sol in gradu 11 Tauri, declinatio hinc septemtrionalis
grad. 15. m. 20. cuius numerus ex tabula benefica 103609
Latitudo loci graduum 38. eius numerus ex tabula eadem
126902 Dria ascensionalis Solis graduum gr. 12. m. 24.
eius sinus 21189 Altitudo Solis esto gr. 5. m. 30. eius sinus
9585. Quibus peractis sic procedo. Multiplico sinum altitudinis
scilicet 9585. in numerum latitudinis loci scilicet 126902
Et produco 1216355670 hinc abiicio quinque ad dextram figuras
Et relinquitur numerus \uncertain{12164} quem rursus multiplico in
numerum declinationis superius habitum scilicet 103609 Et
produco 1260299876 Unde abiectis quinque dextrorsum
figuris restat numerus 12604 quem voco argumentum distantiæ,
Unde quoniam sol caret declinatione, non opus est secunda
multiplicatione: Nam primum productum abiectis quinque
figuris residuat tunc sinum secundum distantiæ Solis a meridiano.
Sed in hoc casu tale argumentum collatum ad sinum differentiæ
ascensionalis indicabit quæsitam distantiam. Nam si æquale
sit ipsi sinui certum est distantiam Solis a meridiano esse
quadrantem. Si argumentum minus sit sinu distantia talis
major est quadrante. Si majus minor. Unde in hoc casu quoniam
argumentum minus est sinu subtraho argumentum a sinu
hoc est subtraho 12603 a sinu prædicto 21189 et supersunt 8586.
qui sinus postulat gr. 4. m. 55. qui iuncti cum quadrante faciunt
\note{« Dria » : abréviation de « Differentia ». Le nombre qui reste après la première division est surchargé : on lit « 12 », un chiffre repassé, puis « 64 » ; « 12164 » est offert parce que $12164 \times 103609 = 1260299876$, le produit écrit à la ligne suivante, et que 1216355670, cinq chiffres ôtés, donne 12163 porté à 12164 par la règle de l'unité ajoutée : c'est une vérification, non une lecture. Le même argument est écrit « 12604 » puis « 12603 » ; les deux sont laissés tels quels. « 1216355670 » est écrit en trois groupes, « 121 6355 670 ». La page s'arrête sur « faciunt ».}

\page{49}
\note{double page ; page de gauche. Elle continue la page de droite de la vue 48 (« qui iuncti cum quadrante faciunt »). La page de droite de la vue 49 est blanche et sans numéro.}
gr. \uncertain{95}. m. 55 distantiam Solis a meridiano quæsitam. Supponatur
altitudo Solis gr. 9. m. 16. cæteris stantibus: huius altitudinis
sinus 16103 qui ductus in 126902 facit 2043502906. Hinc
abiectis quinque figuris restat numerus 20435. qui multiplicatus
in 103609. facit 2117149915. Unde abiectis quinque
figuris restat 21171. qui numerus quoniam æqualis est fere
sinui differentiæ ascensionalis 21189 illinc concludo distantiam
quæsitam a meridiano esse quadrantem.
\note{« gr. 95 » : le premier chiffre est un 9, le second a la forme du 5 de « m. 55 » ; la somme annoncée à la vue 48 (le quadrant et gr. 4 m. 55) ferait gr. 94. m. 55. « 2117149915 » et « 21171 » sont écrits ainsi ; le produit de 20435 par 103609 est 2117249915. Laissés tels quels.}

Item supponatur altitudo Solis gr. 27. m. 50. Cætera ut
prius. Sinus talis altitudinis erit 46690 qui ductus in
126902 facit 5925054380 de quo abiectis ad dextram
quinque figuris, relinquitur numerus 59251 qui ductus
in 103609 producit 6138936859. Unde ablatis quinque
figuris superest 61389 qui quoniam major est sinu differentiæ
ascensionalis 21189 ideo hic ab illo subtrahendo
superstabat g. 66. m. 18. distantia Solis a meridiano quæsita.
Hæc pro declinatione septemtrionali Tandem supponatur
declinatio Solis gr. 15. m. 20. meridionalis cæteris ut in
postremo casu stantibus, ac similiter peractis: Sed tunc
argumentum distantiæ scilicet 61389 iungendum est
cum sinu driæ 21189 et fiet numerus 82578. cuius sinui
debetur arcus gr. 55. m. 40. qui de quadrante sublatus
relinquit tunc gr. 34. m. 20. distantiam Solis a meridiano
quæsitam. Notandum tamen si hæc operatio refertur
ad stellam cuius parallelus tangat solummodo horizontem
cum scilicet complementum eius declinationis æquiparat
latitudinem regionis tunc pro sinu driæ ascensionalis
capiendus est sinus totus: quandoquidem talis stellæ dria
est quadrans. Si autem stellæ parallelus feratur ultra
horizontem tunc similiter faciendum est, sed et sinus
minimæ altitudinis stellæ in meridiano auferendus est
a sinu propositæ altitudinis: Et cum residuo agendum tanquam
sinu altitudinis proposito. Denique pro recto horizonte ubi loci
latitudo nulla est multiplicandus est sinus altitudinis
in numerum tabulæ beneficæ per declinationem Solis acceptum
Et a producto abiicienda more consueto quinque figuræ.
Sic enim relicto sinui debitus arcus de quadrante abiectus
relinquet Solis a meridiano distantiam. Quando autem sol
ibi est in æquatore, quoniam æquator est tunc circulus altitudinis
Iam altitudo a quadrante subtracta residuabit eius
a meridiano distantiam. Memento tamen si abiectæ figuræ
plusquam 50000 denotaverint adiiciendam esse reliquo numero
semper unitatem, quandoquidem tunc quod abiicitur excedit
integri dimidium.
\note{« driæ », « dria » : abréviations de « differentiæ », « differentia », avec un trait. « superstabat » est surchargé. Dans le second exemple, « 61389 » s'obtient de 6138936859 en ôtant cinq chiffres, et « g. 66. m. 18. » est donné sans le calcul intermédiaire. Dans la marge droite, vers le milieu de la page, le cachet rond de la Bibliothèque impériale, « Bibliothèque impériale MSS ».}

Compendium Mathematicæ. omittitur quia
a P. M. Mersenno editum in Apologia scientiarum cap. ultimo.
\note{la ligne « Compendium Mathematicæ » est centrée comme un titre ; « omittitur quia … » est ajouté à sa suite, de la même main et de la même encre. La page s'arrête là ; le reste est blanc.}

\end{document}
